% Windows install and clone steps

\section{Install, once}

The launcher is a Python script run from a Unix-style shell; on Windows that shell is Git Bash,
which comes with Git. Every command in this manual is typed into a Git Bash window.

\begin{enumerate}
  \item \textbf{WSL 2.} Open PowerShell as administrator and run \texttt{wsl --install}. Reboot
  when it asks. Docker Desktop runs its Linux containers inside WSL~2.
  \item \textbf{Docker Desktop.} Download from \url{https://docker.com} and install. In Settings,
  General, keep \emph{Use the WSL 2 based engine} ticked. Start it and wait until the whale icon in
  the tray stops animating; leave it running whenever you use the simulator, or \texttt{docker}
  commands fail with ``cannot connect to the Docker daemon''. Tick \emph{Start Docker Desktop when
  you sign in} in its settings. Check in Git Bash later with \texttt{docker run --rm hello-world}.
  \item \textbf{Git for Windows.} Download from \url{https://git-scm.com} and install with the
  defaults. It provides Git Bash.
  \item \textbf{Python 3.} Install from \url{https://python.org}; on the first installer page tick
  \emph{Add python.exe to PATH}. In Git Bash, \texttt{python --version} should answer.
  \item \textbf{GitHub CLI.} In PowerShell: \texttt{winget install GitHub.cli}. Then, in a new Git
  Bash window, \texttt{gh auth login} and follow the prompts.
  \item \textbf{Foxglove.} Download the installer from \url{https://foxglove.dev/download}, open it
  once, sign in with a free account.
\end{enumerate}

The simulator is a native Windows build and needs a GPU with current drivers. On a laptop with
two GPUs, add the app to Windows \emph{Graphics settings} and set it to the high-performance
one; the app is \texttt{FEB Simulator.exe} under
\texttt{\%USERPROFILE\%\textbackslash.feb-sim\textbackslash app\textbackslash windows}.

\section{Clone and set up}

In Git Bash:

\begin{lstlisting}
git clone https://github.com/Pranman1/feb-racing
cd feb-racing
./feb-sim setup --team "Your Team"
\end{lstlisting}

Pulls the devkit image and downloads the Windows app into \texttt{\textasciitilde/.feb-sim/app/}
(that is \texttt{C:\textbackslash Users\textbackslash you\textbackslash.feb-sim}). A few minutes.
The app is not signed, so the first launch shows a SmartScreen warning: click \emph{More info},
then \emph{Run anyway}, once.

